\documentclass[11pt,a4paper]{article}

\usepackage[T1]{fontenc}
\usepackage[utf8]{inputenc}
\usepackage{lmodern}
\usepackage[a4paper,margin=2.5cm]{geometry}
\usepackage{enumitem}
\usepackage{hyperref}
\usepackage{titlesec}

\setlist[itemize]{noitemsep, topsep=3pt}
\setlist[enumerate]{noitemsep, topsep=3pt}

\title{Chapter 16 --- Vasodilators}
\author{}
\date{}

\begin{document}

\maketitle

\noindent \textbf{Source:} Peck \& Harris, \textit{Pharmacology for Anaesthesia and Intensive Care}, 5th ed., Chapter 16: ``Vasodilators''.

\tableofcontents

\section{Overview}

Vasodilators reduce vascular smooth muscle tone and are used to alter preload, afterload, coronary perfusion, pulmonary vascular resistance, or regional vascular tone. Their effects depend on whether they act mainly on:
\begin{itemize}
    \item \textbf{capacitance vessels (veins)} $\rightarrow$ preload falls
    \item \textbf{resistance vessels (arterioles)} $\rightarrow$ afterload falls
    \item \textbf{both} $\rightarrow$ combined reduction in preload and afterload
\end{itemize}

The chapter groups the drugs broadly into:
\begin{itemize}
    \item nitric oxide donors / nitrates
    \item potassium channel activators
    \item calcium channel antagonists
    \item direct vasodilators
    \item other vasodilators used in specialist settings
\end{itemize}

\section{Sodium nitroprusside}

\subsection{Definition}

Sodium nitroprusside is a potent, short-acting intravenous vasodilator that dilates \textbf{both arteries and veins}.

\subsection{Mechanism}

It releases \textbf{nitric oxide}, which activates \textbf{guanylate cyclase}, increasing intracellular \textbf{cGMP}. This reduces available intracellular calcium in vascular smooth muscle and produces relaxation.

\subsection{Main effects}

\subsubsection*{Cardiovascular}
\begin{itemize}
    \item dilates arteries and veins
    \item reduces systemic vascular resistance
    \item reduces blood pressure
    \item reduces preload
    \item may maintain cardiac output through reflex tachycardia
    \item in heart failure, reduced pre- and afterload may improve output without needing tachycardia
    \item lowers myocardial wall tension and oxygen consumption
    \item no direct positive inotropic effect
    \item tachyphylaxis may occur
\end{itemize}

\subsubsection*{Respiratory}
\begin{itemize}
    \item may inhibit hypoxic pulmonary vasoconstriction
    \item may increase shunt
\end{itemize}

\subsubsection*{CNS}
\begin{itemize}
    \item cerebral vasodilation may raise intracranial pressure
\end{itemize}

\subsection{Kinetics}
\begin{itemize}
    \item given by infusion
    \item extremely short duration of action, usually under 10 minutes
\end{itemize}

\subsection{Toxicity}

Its major limitation is \textbf{cyanide toxicity}.

Nitroprusside metabolism generates:
\begin{itemize}
    \item nitric oxide
    \item cyanide ions
    \item methaemoglobin
\end{itemize}

Cyanide is normally converted to \textbf{thiocyanate} by rhodanase using sulfhydryl donors. Toxicity becomes more likely with prolonged or high-dose infusion, especially when detoxification is overwhelmed.

\subsubsection*{Features suggesting toxicity}
\begin{itemize}
    \item metabolic acidosis
    \item impaired tissue oxygen utilisation
    \item raised venous oxygen content
\end{itemize}

\subsubsection*{Management}
\begin{itemize}
    \item stop infusion
    \item optimise oxygen delivery
    \item antidotal/supportive treatments described in the chapter:
    \begin{itemize}
        \item dicobalt edetate
        \item sodium thiosulfate
        \item nitrites
    \end{itemize}
\end{itemize}

\subsection{Practical exam summary}

\textbf{Fastest controllable mixed arterial and venous vasodilator; useful for rapid BP reduction but limited by cyanide/thiocyanate toxicity and raised ICP.}

\section{Nitrates}

\subsection{Glyceryl trinitrate (GTN)}

\subsubsection*{Definition}

GTN is an \textbf{organic nitrate} used in angina, left ventricular failure, perioperative blood pressure control, and related states.

\subsubsection*{Presentation}

The chapter describes multiple formulations:
\begin{itemize}
    \item sublingual spray
    \item sublingual tablets
    \item buccal / modified-release forms
    \item oral modified-release tablets
    \item transdermal patch
    \item intravenous infusion
\end{itemize}

A practical point highlighted in the chapter is that \textbf{GTN is absorbed by PVC}, so polyethylene administration sets are preferred for infusion.

\subsubsection*{Mechanism}

GTN generates \textbf{nitric oxide}, activating \textbf{guanylate cyclase} and increasing \textbf{cGMP}. The result is reduced intracellular calcium and vascular smooth muscle relaxation.

\subsubsection*{Main effects}

\paragraph{Cardiovascular}
GTN acts \textbf{predominantly on veins}, so it is mainly a \textbf{venodilator}.

Consequences:
\begin{itemize}
    \item reduced preload
    \item reduced venous return
    \item reduced ventricular end-diastolic pressure
    \item reduced ventricular wall tension
    \item reduced myocardial oxygen demand
\end{itemize}

It also:
\begin{itemize}
    \item improves coronary perfusion, especially to subendocardial regions
    \item is useful in myocardial ischaemia and LV failure
    \item causes some arterial dilation, though less prominently than venodilation
\end{itemize}

Potential disadvantages:
\begin{itemize}
    \item hypotension
    \item postural hypotension
    \item reflex tachycardia
    \item excessive hypotension may reduce coronary perfusion
\end{itemize}

\paragraph{CNS}
\begin{itemize}
    \item headache is common
    \item may raise ICP due to cerebral vasodilation
\end{itemize}

\paragraph{GI smooth muscle}
\begin{itemize}
    \item relaxes GI sphincters including the sphincter of Oddi
\end{itemize}

\paragraph{Blood}
\begin{itemize}
    \item methaemoglobinaemia is rare
\end{itemize}

\subsubsection*{Kinetics}
\begin{itemize}
    \item rapidly absorbed sublingually
    \item marked first-pass metabolism limits oral bioavailability
    \item hepatic nitrate reductase contributes to metabolism
    \item tissue thiols are important in bioactivation
\end{itemize}

\subsubsection*{Tolerance}
A key chapter point:
\begin{itemize}
    \item tolerance may develop within 48 hours
    \item likely related to depletion of sulfhydryl groups
    \item a nitrate-free interval reduces this problem
\end{itemize}

\subsubsection*{Practical exam summary}

\textbf{GTN is mainly a venodilator: think preload reduction, anti-anginal action, and usefulness in myocardial ischaemia / LV failure.}

\subsection{Isosorbide dinitrate / mononitrate}

These are longer-acting nitrate preparations used mainly for ongoing anti-anginal treatment rather than rapid titration. The chapter notes substantial first-pass metabolism of isosorbide dinitrate to active mononitrate metabolites.

\section{Potassium channel activator}

\subsection{Nicorandil}

\subsubsection*{Definition}

Nicorandil is a \textbf{potassium channel activator} with an additional \textbf{nitrate moiety}.

\subsubsection*{Uses}
\begin{itemize}
    \item treatment and prophylaxis of angina
    \item also described in the chapter in heart failure and hypertension
\end{itemize}

\subsubsection*{Mechanism}

It works through two linked mechanisms:

\paragraph{1. ATP-sensitive potassium channel activation}
Opening KATP channels causes potassium efflux, membrane hyperpolarisation, closure of calcium channels, and reduced calcium entry.

\paragraph{2. Nitrate-like action}
Its nitrate component stimulates guanylate cyclase, increasing cGMP and relaxing venous capacitance vessels.

\subsubsection*{Main effects}

\paragraph{Cardiovascular}
\begin{itemize}
    \item causes arterial and venous dilation
    \item reduces both preload and afterload
    \item lowers blood pressure
    \item lowers LV end-diastolic pressure
    \item improves normal and collateral coronary blood flow
    \item causes coronary vasodilation without a steal phenomenon
    \item may increase cardiac output in ischaemic heart disease and cardiac failure
    \item useful against torsades de pointes associated with long QT
    \item does not significantly affect AV conduction or myocardial contractility
    \item unlike nitrates, it is not associated with tolerance during prolonged use
\end{itemize}

\paragraph{Other effects}
\begin{itemize}
    \item headache
    \item no adverse effect on lipid profile or glucose control
    \item inhibits ADP-induced platelet aggregation in vitro
    \item may cause giant oral aphthous ulcers
\end{itemize}

\subsubsection*{Practical exam summary}

\textbf{A hybrid anti-anginal drug: KATP opener plus nitrate, giving mixed arterial and venous dilation without nitrate tolerance.}

\section{Calcium channel antagonists}

\subsection{General principles}

The chapter treats calcium channel antagonists as vasodilators because they block \textbf{L-type calcium channels} in the cardiovascular system.

Important distinctions made in the chapter:
\begin{itemize}
    \item L-type channels are important in the plateau phase of the cardiac action potential and in vascular smooth muscle contraction
    \item different drugs have differing affinity for:
    \begin{itemize}
        \item myocardium
        \item SA and AV nodes
        \item vascular smooth muscle
    \end{itemize}
\end{itemize}

This explains why some are mainly:
\begin{itemize}
    \item \textbf{cardiac / nodal depressants}
\end{itemize}
while others are mainly:
\begin{itemize}
    \item \textbf{peripheral and coronary vasodilators}
\end{itemize}

The chapter classifies them chemically as:
\begin{itemize}
    \item \textbf{Class I phenylalkylamines} --- verapamil
    \item \textbf{Class II dihydropyridines} --- nifedipine, amlodipine, nimodipine
    \item \textbf{Class III benzothiazepines} --- diltiazem
\end{itemize}

\subsection{Verapamil}

\subsubsection*{Definition}

Verapamil is a \textbf{racemic mixture} and a synthetic derivative of papaverine.

\subsubsection*{Uses}
\begin{itemize}
    \item certain supraventricular arrhythmias
    \item angina
    \item has been used for hypertension, but its negative inotropic action limits usefulness there
    \item the chapter specifically notes it should \textbf{not} be used in atrial fibrillation complicating Wolff--Parkinson--White syndrome
\end{itemize}

\subsubsection*{Mechanism}

The chapter emphasises stereospecificity:
\begin{itemize}
    \item \textbf{L-isomer}: calcium channel blocking effect, especially at the SA and AV nodes
    \item \textbf{D-isomer}: action on fast sodium channels, giving some local anaesthetic activity
\end{itemize}

\subsubsection*{Main effects}

\paragraph{Cardiovascular}
Verapamil is the most \textbf{cardiac-selective} of the calcium antagonists discussed here.

It:
\begin{itemize}
    \item slows SA and AV nodal conduction
    \item reduces heart rate
    \item prolongs AV conduction time
    \item has a negative inotropic effect
    \item causes some peripheral vasodilation
    \item is a mild coronary vasodilator
    \item lowers blood pressure
\end{itemize}

Risks:
\begin{itemize}
    \item heart block
    \item cardiac failure in patients with impaired ventricular function
    \item ventricular fibrillation in WPW-associated AF
\end{itemize}

\paragraph{CNS}
\begin{itemize}
    \item vasodilates the cerebral circulation
\end{itemize}

\paragraph{Other}
\begin{itemize}
    \item chronic administration may potentiate both depolarising and non-depolarising neuromuscular blockers
\end{itemize}

\subsubsection*{Kinetics}
\begin{itemize}
    \item well absorbed orally
    \item extensive first-pass metabolism reduces oral bioavailability to about 20\%
    \item metabolised to norverapamil, which retains anti-arrhythmic activity
    \item about 90\% protein bound
    \item mainly excreted in urine after metabolism, though some is excreted in bile
\end{itemize}

\subsubsection*{Practical exam summary}

\textbf{Verapamil is primarily a nodal and myocardial calcium antagonist: slows AV conduction, reduces heart rate, depresses contractility, and only secondarily vasodilates.}

\subsection{Nifedipine}

\subsubsection*{Definition}

Nifedipine is a \textbf{dihydropyridine} calcium channel antagonist.

\subsubsection*{Uses}
\begin{itemize}
    \item prophylaxis and treatment of angina
    \item hypertension
    \item Raynaud's syndrome
\end{itemize}

\subsubsection*{Presentation}
\begin{itemize}
    \item capsules
    \item tablets, including sustained-release preparations
    \item the chapter notes that capsule contents may be given sublingually, with rapid onset
\end{itemize}

\subsubsection*{Main effects}

\paragraph{Cardiovascular}
Nifedipine is much more \textbf{vascular-selective} than verapamil.

It:
\begin{itemize}
    \item reduces tone in peripheral and coronary arteries
    \item reduces systemic vascular resistance
    \item lowers blood pressure
    \item increases coronary artery blood flow
    \item produces reflex tachycardia
    \item may increase contractility reflexly
    \item increases cardiac output
\end{itemize}

The chapter notes that these reflex changes may occasionally \textbf{worsen myocardial oxygen supply-demand balance}.

\subsubsection*{Kinetics}
\begin{itemize}
    \item well absorbed orally
    \item first-pass metabolism reduces oral bioavailability to about 60\%
    \item about 95\% protein bound
    \item elimination half-life around 5 hours after oral dosing
    \item mainly excreted in urine as inactive metabolites
\end{itemize}

\subsubsection*{Practical exam summary}

\textbf{Nifedipine is the chapter's prototype vascular calcium antagonist: marked arterial vasodilation, reduced SVR, reflex tachycardia, increased coronary flow.}

\subsection{Nimodipine}

\subsubsection*{Definition}

Nimodipine is a \textbf{more lipid-soluble analogue of nifedipine}.

\subsubsection*{Distinguishing feature}

Because it is more lipid soluble, it can \textbf{cross the blood--brain barrier}.

\subsubsection*{Uses}
\begin{itemize}
    \item prevention and treatment of cerebral vasospasm after subarachnoid haemorrhage
    \item migraine
    \item may be given orally or intravenously
\end{itemize}

\subsubsection*{Mechanistic emphasis in the chapter}

The chapter suggests its effect may depend not only on vasodilation but also on blocking a \textbf{calcium-dependent cascade of cellular injury}, thereby limiting cell damage.

\subsubsection*{Practical exam summary}

\textbf{Nimodipine is the CNS-directed dihydropyridine: think cerebral vasospasm after SAH because it reaches the brain.}

\subsection{Diltiazem}

\subsubsection*{Definition}

Diltiazem is a \textbf{benzothiazepine} calcium antagonist with effects intermediate between verapamil and nifedipine.

\subsubsection*{Uses}
\begin{itemize}
    \item prophylaxis and treatment of angina
    \item hypertension
\end{itemize}

\subsubsection*{Main effects}

\paragraph{Cardiovascular}
Diltiazem acts in both the \textbf{heart} and the \textbf{peripheral circulation}.

It:
\begin{itemize}
    \item prolongs AV conduction time
    \item reduces contractility, but less than verapamil
    \item reduces systemic vascular resistance
    \item lowers blood pressure
    \item usually does not cause marked reflex tachycardia
    \item increases coronary blood flow by coronary vasodilation
\end{itemize}

\subsubsection*{Kinetics}
\begin{itemize}
    \item almost completely absorbed orally
    \item first-pass metabolism reduces bioavailability to about 50\%
    \item produces an active metabolite, desacetyldiltiazem
    \item about 75\% protein bound
    \item metabolite is excreted in urine
    \item the chapter notes both hepatic metabolism and renal elimination
\end{itemize}

\subsubsection*{Practical exam summary}

\textbf{Diltiazem sits between verapamil and nifedipine: some AV nodal slowing and negative inotropy, but also useful vasodilation with less reflex tachycardia.}

\subsection{Comparative haemodynamic pattern from the chapter}

The chapter's table is worth converting into words:

\subsubsection*{Verapamil}
\begin{itemize}
    \item blood pressure: $\downarrow$
    \item heart rate: $\downarrow$
    \item AV conduction time: $\uparrow$
    \item myocardial contractility: $\downarrow\downarrow$
    \item peripheral/coronary vasodilation: mild
\end{itemize}

\subsubsection*{Nifedipine}
\begin{itemize}
    \item blood pressure: $\downarrow$
    \item heart rate: unchanged or $\uparrow$
    \item AV conduction time: unchanged or $\uparrow$ slightly
    \item myocardial contractility: unchanged or $\uparrow$ reflexly
    \item peripheral/coronary vasodilation: marked
\end{itemize}

\subsubsection*{Diltiazem}
\begin{itemize}
    \item blood pressure: $\downarrow$
    \item heart rate: $\downarrow$
    \item AV conduction time: $\uparrow$
    \item myocardial contractility: $\downarrow$
    \item peripheral/coronary vasodilation: moderate
\end{itemize}

\subsubsection*{Nimodipine}
The chapter does not provide a separate haemodynamic table entry, but conceptually it behaves as a \textbf{dihydropyridine}, with the distinguishing feature being \textbf{CNS penetration} and use in cerebral vasospasm.

\section{Direct vasodilators}

\subsection{Hydralazine}

\subsubsection*{Uses}
\begin{itemize}
    \item chronic hypertension
    \item severe chronic heart failure with other agents
    \item intravenous treatment of acute hypertension in pre-eclampsia
\end{itemize}

\subsubsection*{Mechanism}

The exact mechanism is stated to be uncertain, but the chapter links it to:
\begin{itemize}
    \item activation of guanylate cyclase
    \item increased cGMP
    \item reduced available intracellular calcium
    \item vasodilation
\end{itemize}

\subsubsection*{Main effects}
\begin{itemize}
    \item mainly arteriolar dilation
    \item fall in SVR
    \item little effect on capacitance vessels
    \item therefore less postural hypotension than venodilators
    \item reflex tachycardia and increased cardiac output are common
    \item $\beta$-blockade may attenuate these responses
\end{itemize}

Other adverse effects:
\begin{itemize}
    \item increased cerebral blood flow
    \item sodium and water retention
    \item nausea and vomiting
    \item peripheral neuropathy
    \item blood dyscrasias
    \item lupus-like syndrome, especially with long-term use
\end{itemize}

\subsubsection*{Practical exam summary}

\textbf{Hydralazine is a direct arteriolar dilator: afterload reduction, reflex tachycardia, fluid retention, lupus-like toxicity.}

\subsection{Minoxidil}

The chapter describes minoxidil as similar in principle to hydralazine.

\subsubsection*{Main points}
\begin{itemize}
    \item potent arteriolar vasodilator
    \item used in severe hypertension
    \item also associated with hair growth
    \item may cause hypertrichosis
    \item unlike hydralazine, does not cause a lupus-like syndrome
\end{itemize}

\subsubsection*{Practical exam summary}

\textbf{Potent arteriolar dilator remembered for hypertrichosis.}

\subsection{Diazoxide}

\subsubsection*{Uses}
\begin{itemize}
    \item hypertensive emergencies, especially with renal disease
    \item intractable hypoglycaemia
    \item malignant islet-cell tumours
\end{itemize}

\subsubsection*{Mechanism}

The chapter attributes its hypotensive effect mainly to altered \textbf{cAMP} in arterioles, with reduced calcium influx possibly contributing.

\subsubsection*{Main effects}

\paragraph{Cardiovascular}
\begin{itemize}
    \item arteriolar vasodilation
    \item little effect on capacitance vessels
    \item reduced BP
    \item increased heart rate and cardiac output
    \item postural hypotension is not prominent
\end{itemize}

\paragraph{Metabolic}
\begin{itemize}
    \item inhibits insulin secretion
    \item raises blood glucose
    \item increases catecholamines, renin, and aldosterone
    \item causes fluid retention
\end{itemize}

\subsubsection*{Practical exam summary}

\textbf{Arteriolar dilator with hyperglycaemic action due to inhibition of insulin release.}

\section{Other vasodilators}

\subsection{Sildenafil}
\begin{itemize}
    \item used in pulmonary arterial hypertension and erectile dysfunction
    \item inhibits PDE-5
    \item increases cGMP
    \item relaxes pulmonary vascular smooth muscle and corpus cavernosum
\end{itemize}

\subsection{Bosentan}
\begin{itemize}
    \item used in pulmonary arterial hypertension
    \item dual endothelin receptor antagonist
    \item reduces pulmonary and systemic vascular resistance
\end{itemize}

\section{High-yield comparisons}

\subsection{Nitroprusside vs GTN}

\subsubsection*{Nitroprusside}
\begin{itemize}
    \item arterial + venous dilator
    \item very rapid and titratable
    \item ideal for immediate BP control
    \item toxicity: cyanide / thiocyanate
\end{itemize}

\subsubsection*{GTN}
\begin{itemize}
    \item mainly venodilator
    \item best for angina and myocardial ischaemia
    \item reduces preload and myocardial oxygen demand
    \item tolerance develops
\end{itemize}

\subsection{Verapamil vs nifedipine vs diltiazem vs nimodipine}

\subsubsection*{Verapamil}
\begin{itemize}
    \item most nodal / cardiac
    \item $\downarrow$ HR
    \item $\uparrow$ AV conduction time
    \item $\downarrow$ contractility
    \item modest vasodilation
\end{itemize}

\subsubsection*{Nifedipine}
\begin{itemize}
    \item most vascular
    \item marked arterial vasodilation
    \item reflex tachycardia
    \item $\uparrow$ coronary flow
\end{itemize}

\subsubsection*{Diltiazem}
\begin{itemize}
    \item intermediate profile
    \item some AV nodal slowing
    \item some negative inotropy
    \item useful vasodilation without much reflex tachycardia
\end{itemize}

\subsubsection*{Nimodipine}
\begin{itemize}
    \item dihydropyridine with high lipid solubility
    \item penetrates CNS
    \item used for cerebral vasospasm after SAH
\end{itemize}

\subsection{Blindspots and additions from other anaesthesia texts}

Compared with the broader anaesthesia books, Peck \& Harris is strongest on \textbf{mechanism}, \textbf{vessel selectivity}, \textbf{core haemodynamic effects}, and \textbf{classic adverse effects}. Its main blindspots are not usually errors of fact, but rather areas of \textbf{clinical under-emphasis}. These are worth making explicit for revision.

\subsubsection*{1. Perioperative management is under-emphasised}
Peck \& Harris is primarily a pharmacology text, so it gives less detail on how these drugs are handled around anaesthesia and surgery.

Important practical additions from the other books include:
\begin{itemize}
    \item nitrates are generally \textbf{continued perioperatively}, using intravenous or transdermal routes if needed
    \item calcium channel blockers are generally \textbf{continued preoperatively} and restarted early postoperatively
    \item \textbf{sublingual nifedipine} is treated cautiously because abrupt haemodynamic change may worsen myocardial ischaemia
\end{itemize}

Thus, Peck \& Harris explains \emph{what the drugs do}, but the other texts explain more clearly \emph{what the anaesthetist should do with them perioperatively}.

\subsubsection*{2. Nimodipine is not fully embedded in neuroanaesthetic practice}
Peck \& Harris correctly highlights that nimodipine is:
\begin{itemize}
    \item a dihydropyridine
    \item highly lipid soluble
    \item able to cross the blood--brain barrier
    \item used in cerebral vasospasm after subarachnoid haemorrhage
\end{itemize}

However, the other anaesthesia texts place this more firmly in the context of \textbf{subarachnoid haemorrhage management}, delayed neurological deficit, and the wider neuroanaesthetic pathway. In other words, Peck \& Harris gives the pharmacology, whereas the other books better supply the \textbf{clinical framework}.

\subsubsection*{3. Anti-anginal class comparison is less developed than in the general anaesthesia texts}
Peck \& Harris describes individual drugs well, especially:
\begin{itemize}
    \item verapamil as predominantly nodal / myocardial
    \item nifedipine as predominantly vascular
    \item diltiazem as intermediate
    \item nimodipine as CNS-directed
\end{itemize}

What the other books add more clearly is the \textbf{anti-anginal comparison between whole classes}, especially:
\begin{itemize}
    \item effects on preload
    \item effects on afterload
    \item effects on coronary blood flow
    \item comparison of calcium channel blockers with nitrates and beta blockers
\end{itemize}

This matters for exams because it improves class-level recall, not just drug-by-drug recall.

\subsubsection*{4. Modern perioperative vasodilators are underrepresented}
A major limitation of Peck \& Harris is that it is weighted toward the \textbf{classic vasodilators}. The other anaesthesia texts give greater emphasis to vasodilators that are highly relevant in contemporary perioperative and critical care practice, especially:
\begin{itemize}
    \item \textbf{nicardipine}
    \item \textbf{clevidipine}
    \item \textbf{inhaled nitric oxide}
\end{itemize}

These are particularly important for:
\begin{itemize}
    \item perioperative hypertension
    \item controlled hypotension
    \item cardiothoracic anaesthesia
    \item pulmonary hypertension and right ventricular dysfunction
\end{itemize}

So Peck \& Harris is excellent for classical pharmacology, but less complete for the \textbf{modern anaesthetic vasodilator toolbox}.

\subsubsection*{5. Pulmonary vasodilation is not as fully integrated}
Peck \& Harris includes \textbf{sildenafil} and \textbf{bosentan}, but the other books develop the pulmonary vasodilator theme more clearly by emphasising:
\begin{itemize}
    \item \textbf{inhaled nitric oxide} as a \textbf{selective pulmonary vasodilator}
    \item its use in reversible pulmonary hypertension
    \item its usefulness where improved perfusion of ventilated lung regions is desirable
\end{itemize}

Therefore, if revising pulmonary vascular manipulation for ICU, cardiothoracic, or one-lung ventilation, Peck \& Harris should be supplemented.

\subsubsection*{6. Disease-specific anaesthetic applications are relatively sparse}
The other books provide more examples of where the pharmacology matters in specific clinical settings. Useful additions include:
\begin{itemize}
    \item \textbf{verapamil} in conditions where negative inotropy and rate control are clinically relevant
    \item \textbf{hydralazine} in the \textbf{denervated transplanted heart}, where reflex chronotropic responses are absent
    \item \textbf{diazoxide} in endocrine settings such as refractory hypoglycaemia or insulinoma
    \item \textbf{nicorandil} placed more clearly within the anti-anginal treatment strategy
\end{itemize}

Peck \& Harris explains the pharmacology, but the wider anaesthesia texts more often explain the \textbf{clinical scenario in which the pharmacology becomes important}.

\subsubsection*{7. Practical administration and drug selection are less detailed}
Peck \& Harris includes some useful practical details, for example:
\begin{itemize}
    \item GTN absorption by PVC
    \item nitrate tolerance
    \item nitroprusside cyanide metabolism
\end{itemize}

However, compared with the handbooks and larger anaesthesia texts, it gives less emphasis to:
\begin{itemize}
    \item real-world titration strategy
    \item why one vasodilator is chosen over another intraoperatively
    \item operative context, such as cardiothoracic surgery, intubation responses, or hypertensive crises
\end{itemize}

Thus, it is stronger on \textbf{pharmacological description} than on \textbf{drug selection in theatre or ICU}.

\subsubsection*{8. Nitroprusside is well covered, but alternatives are not emphasised enough}
Peck \& Harris gives a strong account of sodium nitroprusside, including:
\begin{itemize}
    \item nitric oxide / cGMP mechanism
    \item mixed arterial and venous dilation
    \item cyanide and thiocyanate toxicity
    \item metabolic acidosis and impaired tissue oxygen utilisation as markers of toxicity
\end{itemize}

The blindspot is not in the nitroprusside material itself, but in the fact that the other anaesthesia texts more clearly highlight that, in practice, alternative agents may be preferred in some settings because of:
\begin{itemize}
    \item easier titration
    \item greater vascular selectivity
    \item fewer toxicity concerns
\end{itemize}

\subsubsection*{9. Overall revision implication}
For examination purposes, Peck \& Harris should be treated as the \textbf{core pharmacology source} for vasodilators because it explains mechanism and haemodynamic logic very well. It is best supplemented by other anaesthesia texts for:
\begin{itemize}
    \item perioperative continuation or withholding of therapy
    \item modern vasodilators used in current practice
    \item pulmonary vasodilator strategy
    \item disease-specific anaesthetic application
    \item practical intraoperative drug choice
\end{itemize}

In summary, the main blindspots are:
\begin{enumerate}
    \item perioperative handling of nitrates and calcium channel blockers
    \item modern agents such as nicardipine, clevidipine, and inhaled nitric oxide
    \item fuller pulmonary vasodilator practice
    \item stronger disease-specific clinical anchoring
    \item practical anaesthetic decision-making
\end{enumerate}

\end{document}
